\section{検知精度への影響の評価計画}

5章は検知条件そのものを比較した．
本章はその差が実際のスキャン結果にどう現れるかを測る計画である．結果は発表時に報告する．

同一の基盤\cite{vuls-data-update,vuls-data-db}を用い，入力データソースのみを変えた3種のデータベースを構築する．
\textbf{DB-nvd} はNVDの \texttt{configurations} のみ，\textbf{DB-vc} はVulnCheck NVD++ の \texttt{configurations} と \texttt{vcConfigurations}，\textbf{DB-primary} は表\ref{tab:sources}の一次データソース（6章の方式で構造化したものを含む）である．

インベントリは2種を用いる．
\textbf{合成インベントリ}は5.2節の境界値法をそのまま用い，一次データソースの影響版と修正版から「検知されるべき」「検知されるべきでない」のラベル付きバージョンを生成する．
5章はこれを検知条件の集合演算で評価したが，本章では実際にスキャナを動かし，データベース構築と照合実装を含む経路全体での未検知・過剰検知を測る．
\textbf{実環境インベントリ}はVulsのスキャン対象として整備しているコンテナ・仮想環境\cite{vuls-targets}であり，運用に近い条件での差分を確認する．

指標は5.4節と同じ未検知率・過剰検知率に加え，スキャン結果の件数差，および過剰検知のうちバージョン指定を持たない条件（\texttt{-}，\texttt{*}）に起因する割合を報告する．
さらに公開年で2023年以前 / 2024年 / 2025年 / 2026年に分割して同じ指標を算出し，「どの時点までのNVDを信頼してよいか」を検知精度の観点から述べ，併せて4.5節のCNA別の差が検知精度の差として現れるかを確認する．
